% Theme: Flashcard  --  best for quizzes: question -> hint -> answer
% Build:  ./build.sh example-flashcard.tex   (or compile on Overleaf)
\documentclass[12pt]{beamer}
\usetheme{igflashcard}

\ighandle{@quantum_danesh}
\igseries{Quiz cuántico}
\igissue{07}

\begin{document}

\coverslide{Quiz 07}
  {¿Es la transpuesta un canal cuántico?}
  {Piensa tu respuesta antes de deslizar.}

\begin{frame}{¿Cuál es correcta?}{Pregunta}
Sea $T(\rho) = \rho^{\top}$ la transpuesta, actuando sobre un qubit.
\begin{enumerate}[(A)]
  \item Sí: es completamente positiva y preserva la traza.
  \item Es positiva y preserva la traza, pero no es completamente positiva.
  \item Ni siquiera es positiva.
  \item Es completamente positiva, pero no preserva la traza.
\end{enumerate}
\end{frame}

\begin{frame}{¿Sigue siendo un estado?}{Pista}
\begin{notebox}
Aplica $\mathrm{id}\otimes T$ a la mitad de
\[
  \ket{\Phi^+} = \tfrac{1}{\sqrt2}\big(\ket{00} + \ket{11}\big)
\]
y mira los autovalores de lo que sale.
\end{notebox}
Si alguno es negativo, el resultado ya no es una matriz densidad.
\end{frame}

\begin{frame}{Es la (B)}{Respuesta}
\begin{enumerate}[(A)]
  \item {\color{igMuted}Sí: es completamente positiva y preserva la traza.}
  \item \hl{Es positiva y preserva la traza, pero no es completamente
        positiva.} {\color{igAccent}\checkmark}
  \item {\color{igMuted}Ni siquiera es positiva.}
  \item {\color{igMuted}Es completamente positiva, pero no preserva la traza.}
\end{enumerate}
\end{frame}

\begin{frame}{La cuenta}{Por qué}
\begin{align*}
  (\mathrm{id}\otimes T)\big(\ketbra{\Phi^+}\big)
  &= \frac12\sum_{i,j}\ket{i}\!\bra{j}\otimes\ket{j}\!\bra{i}\\
  &= \frac12\,\mathrm{SWAP}
\end{align*}
$\mathrm{SWAP}$ tiene autovalores $+1$ (tres veces) y $-1$ (una vez):
\keyeq{\operatorname{spec} = \big\{\tfrac12,\ \tfrac12,\ \tfrac12,\ \mhl{-\tfrac12}\big\}}
Un autovalor negativo: \key{no es completamente positiva}, aunque
$\rho^{\top}$ conserve espectro y traza.
\end{frame}

\begin{frame}{Positivo no es completamente positivo}{Para recordar}
Por eso la transpuesta no es un canal, aunque mande estados a estados.
\begin{thmbox}[Criterio de Peres--Horodecki]
Si la traspuesta parcial de un estado tiene un autovalor negativo, el
estado está entrelazado. Para dos qubits, también vale al revés.
\end{thmbox}
\end{frame}

\closingslide{¿Acertaste?}
  {Cuéntame en los comentarios qué opción elegiste antes de deslizar.}

\end{document}
